%!TEX root = ../thesis_main.tex
%!TEX encoding = UTF-8 Unicode

\section{Equilibrium model description}
\label{sec:eqmodel}

\gls{PCE} \ie A thermodynamic equilibrium model is used to predict the impact of steam and oxygen injection on the gasification process, and to support the interpretation of experimental results. In this model, the composition of syngas is calculated assuming that all chemical reactions proceed to their equilibrium, disregarding chemical kinetics and process dynamics.  \cite{Evans2012} \citet{Evans2012}

The stoichiometric approach is used for the equilibrium model. The model reactants are: biomass on dry and ash-free basis ($\mathit{daf}$) represented by the formula \ce{CH_yO_xN_z}, biomass moisture \ce{H2O_l}, and the gasification agents \ce{N2}, \ce{O2} and steam \ce{H2O_g}. At the gasifier outlet, raw syngas is assumed to contain only \ce{CO}, \ce{CO2},  \ce{CH4}, \ce{H2}, \ce{H2O_g} and \ce{N2}, and residual solid carbon \ce{C_s} can subsist in the products if the gasification is incomplete. The 6 unknown coefficients of the products composition (Eq. \ref{eq:main_eq_model}) are solved using the elemental mass balances of \ce{C}, \ce{H} and \ce{O}, combined with the equilibrium constants of three reactions. The equilibrium reactions considered are Boudouard reaction (Eq. \ref{eq:Boudouard}), the Water-Gas Shift (WGS) (Eq. \ref{eq:WGS}) and steam methane reforming (Eq. \ref{eq:SMR}). The equilibrium constants are calculated using NIST-JANAF tables \cite{Chase1998}, for a given equilibrium temperature.

\begin{multline}
\label{eq:main_eq_model} \ce{CH_yO_xN_z} + m \ce{H2O_l} + w\left(\ce{O2} + \xi\ce{N2}\right) + s \ce{H2O_g} \\
\Rightarrow a_1\mathrm{CO} + a_2\ce{CO2} + a_4\ce{CH4} + b_1\ce{H2} + b_2\ce{H2O_g} + \left(\frac{z}{2}+\xi w\right)\ce{N2} + (a_0+u_{\ce{C}})\ce{C_s}
\end{multline}
\begin{eqnarray}
\label{eq:Boudouard}\ce{C_s + CO2 & <=>& 2CO}\\
\label{eq:WGS}\ce{CO + H2O &<=> &CO2 + H2}\\
\label{eq:SMR}\ce{CH4 + H2O &<=> &CO + 3H2}
\end{eqnarray}

The equilibrium temperature is determined, iteratively with the composition of the products, using the energy conservation between the reactants $R$ and the products $P$ (Eq. \ref{eq:energy_conservation}). In this equation, $h^{\circ}_i(T)$ [kJ/kmol] is the total enthalpy (formation and sensible) of species $i$ at temperature $T$ [K] and $n_i$ [kmol/kmol$_{\ce{CH_yO_xN_z}}$] is the amount of species $i$ in the reactants or products. $H_\text{loss}$ [kJ/kmol$_{\ce{CH_yO_xN_z}}$] is an exogenous heat loss, $A$ [kg$_\text{ash}$/kmol$_{\ce{CH_yO_xN_z}}$] is the biomass ash content and $h_A(T_{eq})$ [kJ/kg] is the sensible enthalpy of ash at temperature $T$ [K], determined using a correlation (Eq. \ref{eq:ash_enthalpy}) proposed by \citet{Song2013}.

\begin{equation}
\label{eq:energy_conservation} \sum_{i\in R} n_i h_{i}^{\circ}(T_{i,0}) = \sum_{i\in P} n_i h_{i}^{\circ}(T_{eq}) + A h_A(T_{eq}) + H_\text{loss}
\end{equation}
\begin{equation}
\label{eq:ash_enthalpy} h_A(T) = 0.0002155  T^2 + 0.7618 T - 254.20
\end{equation}

Steam injection is defined by the Steam-to-Biomass mass ratio (SBR) (Eq. \ref{eq:SBR}), which determines the coefficient $s$, and by the steam temperature $T_{st}$. Oxygen injection is represented by an equivalent air enrichment affecting the \ce{N2}/\ce{O2} ratio, $\xi$. The amount of oxygen $w$ is defined by the air-fuel Equivalence Ratio (ER) (Eq. \ref{eq:ER}), which is either imposed or computed to maximize the Cold Gas Efficiency (CGE) (Eq. \ref{eq:CGE}). Other inputs to the model are the biomass characteristics (elemental composition, $\mathrm{LHV}$, moisture and ash content), the process pressure $p$ (1\,bar), the inlet temperature of air and biomass, and the exogenous heat losses $H_\text{loss}$.

\begin{eqnarray}
\label{eq:SBR}\mathrm{SBR}&=&\frac{\dot m_{\ce{H2O_g}}}{\dot m_{\ce{CH_yO_xN_z}}}\quad = \quad s\frac{M_{\ce{H2O}}}{M_{\ce{CH_yO_xN_z}}}\\
\label{eq:ER}\mathrm{ER}&=&\frac{1}{1 + \frac y4 - \frac x2}\frac{\dot n_{\ce{O2}}}{\dot n_{\ce{CH_yO_xN_z}}}\;\; =\;\;\frac{w}{1 + \frac y4 - \frac x2}\\
\label{eq:CGE} \mathrm{CGE}&=&\frac{P_\text{syngas}}{P_\text{biomass,\textit{wb}}}\quad = \quad\frac{a_1\mathrm{LHV}_{\ce{CO}} + b_1\mathrm{LHV}_{\ce{H2}} + a_4\mathrm{LHV}_{\ce{CH4}}}
{\mathrm{LHV}_{\ce{CH_yO_xN_z}}-mh_{lv,\ce{H2O}}}
\end{eqnarray}

Two parameters can be adapted to obtain results more representative of real conditions: the exogenous heat losses $H_\text{loss}$ and the minimum amount of unconverted carbon $u_{\ce{C}}$. The amount of unconverted carbon is set to $u_{\ce{C}}=8\%$, based on typical experimental values. This corresponds to the fraction of the carbon that is usually found in the ashes: as char progresses downwards in the reduction zone, it moves towards a lower temperature zone, preventing the complete conversion of carbon. Regarding the heat losses, the best match between the model and the experimental observations is found by setting $H_\text{loss}=2.42$\,MJ/kg$_\text{wood,\textit{daf}}$.

%% Remarque: la médiane est en fait à 9% mais personne ne va vérifier + le graphe indique plutôt 8%

%Heat losses are mainly due to the wall heat transfer by conduction. Based on the gasifier wall layers thickness and individual conductivity, and typical internal and external temperatures, these losses are estimated to 6\,kW. However, the best match between the model and the experimental observations is found by setting $H_\text{loss}=2.42$\,MJ/kg$_\text{wood,\textit{daf}}$, corresponding to 13\% of the input biomass LHV, and to 30\,kW at 50\,kg/h.

In the results described in the next section, the Equivalence Ratio (ER) is optimized for each value of the equivalent air enrichment (from air to \ce{O2}) without steam. When adding steam, the ER is kept constant to the situation at SBR=0, neglecting any impact of steam on the optimal ER. If the ER is left to be optimized, it rapidly decreases when increasing the SBR, as the added steam theoretically replaces \ce{O2} as gasifying agent. Steam-involved reactions are mainly endothermic, and the equilibrium temperature drops if the model is allowed to seek such optimum: the model then predicts a very high production of \ce{CH4} which is not realistic as the assumption that reactions proceed close to the equilibrium becomes unacceptable. To avoid this effect, the ER was forced not to vary with the SBR.


\section{Gasification with varying gasification agents}

\subsection{Steam and oxygen}


\subsection{Hydrogen and \ce{CO2}}








